\section{The exponent of matrix multiplication}\label{sec:intro}

The product of two square matrices of side \(n\) has \(n^2\) entries.
Computing the dot product for each entry separately takes order \(n^3\)
arithmetic operations. Strassen's seven-product identity showed that
sharing work among the entries can save a power of \(n\). How far can
those savings go? The conjecture \(\omega=2\) asks whether, for every
\(\varepsilon>0\), the product can be computed in \(O(n^{2+\varepsilon})\)
operations. We show that the arithmetic exponent exceeds the exponent of the output
size by exactly one quarter.

\begin{theorem}\label{thm:square}
Over every field \(F\), the exponent of square matrix multiplication is
\(\omega_F=9/4\).
\end{theorem}

Here \(\omega_F\) is the infimum of the real numbers \(s\) for which two \(n\times n\)
matrices can be multiplied in \(O_{F,s}(n^s)\) scalar operations. We use
formal arithmetic circuits. Their inputs are independent variables,
their constants lie in \(F\), and their gates can add, subtract,
multiply or divide by a nonzero rational function. Each output must
agree with the corresponding entry of the product as an identity of
rational functions. This convention applies also over finite fields.

The \(9/4\) upper bound was established over \(\mathbb C\) by OpenAI
\cite{nine} and over every field by Amanov, van den Berg and Zuiddam
\cite{abvz}, following L\"u's extension to arbitrary characteristic
\cite{lu}. We prove the matching lower bound. In fact, for every
\(n\ge2\) a circuit for two \(n\times n\) matrices requires at least
\(\lceil n^{9/4}/2\rceil\) multiplication or division gates, even if
linear operations are free. The upper bound gives circuits of size
\(O_{F,\varepsilon}(n^{9/4+\varepsilon})\) for each \(\varepsilon>0\).

\par\Needspace{8\baselineskip}
\subsection{Rank at each matrix size}

A bilinear algorithm forms products of linear combinations of entries
from the first matrix with linear combinations from the second. It
then combines these products linearly to give all the output entries.
The least number of products is the rank of the matrix multiplication
tensor
\[
             T_n=\sum_{i,j,k=1}^{n}
                    e_{ij}\otimes e_{jk}\otimes e_{ki}.
\]
The third factor pairs the output at position \((i,k)\) with \(e_{ki}\).
Grouping corresponding factors gives \(T_n^{\otimes j}\simeq T_{n^j}\).
Consequently an improvement at one fixed matrix size can be reused at
arbitrarily large sizes.

Geometric border rank allows the tensor to belong to the Zariski closure,
over \(\overline F\), of tensors computed with a given number of products.
Asymptotic rank over \(F\) is the limit of the \(j\)-th roots of the
ranks over \(F\) of the \(j\)-th tensor powers.

\begin{corollary}\label{cor:tight}
For every \(n\ge2\) and every field \(F\), the geometric border rank of
\(T_n\) is at least \(n^{9/4}\), and the asymptotic rank of \(T_n\) over
\(F\) equals \(n^{9/4}\). Every formal arithmetic circuit over \(F\)
for this product has at least \(\lceil n^{9/4}/2\rceil\) nonscalar
multiplication or division gates.
\end{corollary}

In particular, an exact bilinear algorithm at side \(n\) uses at least
\(\lceil n^{9/4}\rceil\) products. The factor of two for general circuits
comes from extracting the bilinear part of a gate. A multiplication
or division introduces at most two new products of linear forms
(Lemma~\ref{lem:bilinearize}).

The square bound also restricts how wide a rectangular product can be
while its cost still has exponent two. Write \(T_{a,b,c}\) for the
multiplication tensor of an \(a\times b\) matrix and a \(b\times c\)
matrix. Let \(\alpha_F\) be the supremum of the \(\beta\ge0\) for which
multiplying an \(n\times\lfloor n^\beta\rfloor\) matrix by a
\(\lfloor n^\beta\rfloor\times n\) matrix has exponent two in \(n\).

\begin{corollary}\label{cor:rectangular}
For all positive integers \(a,b,c\), the geometric border rank of
\(T_{a,b,c}\) and its asymptotic rank over \(F\) are at least
\[
              \max\{ab,bc,ca,(abc)^{3/4}\}.
\]
Over every field, \(1/2\le\alpha_F\le2/3\).
\end{corollary}

Three copies of the rectangular tensor, with the roles of the factors
cycled, multiply to the square tensor \(T_{abc}\). This explains the
power \(3/4\) in Corollary~\ref{cor:rectangular}. The lower endpoint for
\(\alpha_F\) comes from the known algorithm for inner dimension
\(\lfloor\sqrt n\rfloor\)~\cite{uznanski,abvz}. Section~\ref{sec:finite-ranks}
proves the two corollaries and explains the consequence for Strassen's
conjecture on tight tensors.

\subsection{Why a faster matrix product would test lists too cheaply}

Suppose one could save a fixed power of \(n\) below \(9/4\). A single
bilinear algorithm witnessing a strict saving could then be specialized
to a finite field. Over that field we compare two estimates for a
different computation. Given a list of \(2^w\) words of length \(2w\)
bits, return one if all the words are distinct and zero if any word repeats. Sorting the words and
checking adjacent pairs already gives a small network for this test.
The question is how much of its cost can be shared among several lists
whose inputs vary independently.

We measure cost by counting the vertices and wires in a tensor network
with at most three bits incident at each vertex. The network sums over
its internal bits and, for each setting of its exterior bits, must give
the exact zero-one answer. With \(t\) lists it must give the product of
their \(t\) answers. Its internal wiring may connect different lists.
Even with that freedom, its volume per input bit exceeds \(w^\sigma\)
for some fixed \(\sigma>0\), uniformly for \(1\le t\le2^w\) and all
large \(w\) (Theorem~\ref{thm:many}). The proof reduces the problem to
routing tagged records. The companion results on tensor geometry and
network complexity control this routing or turn cheap routing into a
violation of a counting lower bound~\cite{oaG,oaP}.

The hypothetical matrix multiplication algorithm gives an estimate in
the opposite direction. Consider a trivalent network on \(e\) exterior bits with size at most
\(e^{1.05}\), where \(e\) is large. We can represent a tensor power of
its table in one new network. Per copy, the new network saves a power of the logarithm on
the old size and charges an arbitrarily small power of the logarithm
per exterior bit (Theorem~\ref{thm:shaving}). The replicated bits remain
individually accessible, so we can apply the construction again. Starting
with the sorting test and compressing a fixed number of times gives
volume \(O(w^\eta)\) per input bit for any chosen \(0<\eta<\sigma\).
The batch still has at most \(2^w\) lists. Section~\ref{sec:obstruction}
gives the short contradiction and the passage back to arbitrary fields.
Most of the proof is devoted to deriving this compression from a strict
saving in matrix multiplication.

\subsection{Three wires and a shared label}

To see how the exponent governs network identities, consider a number
assigned to each tensor which is multiplicative under tensor products,
additive under direct sums and monotone under linear maps on the
factors. With the usual normalization this is a spectral point in
Strassen's sense~\cite{spectrum}. Its logarithm is a \emph{price}.
A wire shares an index between two factors. Its price per bit defines
a distance between their sites. The companion paper~\cite{oaG}
extends these distances to spaces of nonpositive curvature.
Section~\ref{sec:prices-background} recalls the properties used here.

Figure~\ref{fig:three-wires} gives two descriptions of \(T_n\). Each of \(i,j,k\) occurs at exactly two factors. The
price of the tensor is therefore the sum of three wire lengths
times \(\log_2 n\). A rank decomposition with \(r\) terms supplies
the same tensor from a single shared label, the choice of term, with
\(r\) possible values. If the price per bit of this three-way label
is \(D\), monotonicity gives
\[
      (\text{sum of the three wire lengths})\,\log_2 n
                                     \ \le\ D\,\log_2 r .
\]
Thus the ratio of a triangle's perimeter to its shared-label rate is a
lower bound on the exponent. Prices coming from Euclidean flats give
ratio two~\cite{oaG}. A different geometry is needed to reach \(9/4\).

\begin{figure}[!htbp]
\centering
\begin{tikzpicture}[x=1cm,y=1cm,font=\small,
   site/.style={circle,draw=ink,fill=white,minimum size=6mm,inner sep=1.5pt},
   conn/.style={thick,draw=sea}]
  \begin{scope}[xshift=-2.75cm]
   \coordinate (A) at (-1.48,0); \coordinate (B) at (1.48,0);
   \coordinate (C) at (0,1.48);
   \draw[conn] (A)--node[below] {\(j\)} (B)
                --node[above right] {\(k\)} (C)
                --node[above left] {\(i\)} (A);
   \node[site] at (A) {\(A\)}; \node[site] at (B) {\(B\)};
   \node[site] at (C) {\(C\)};
   \node at (0,-.85) {Three indices of \(T_n\)};
  \end{scope}
  \begin{scope}[xshift=2.75cm]
   \coordinate (A) at (-1.48,0); \coordinate (B) at (1.48,0);
   \coordinate (C) at (0,1.48); \coordinate (O) at (0,.52);
   \draw[conn] (O)--(A) (O)--(B) (O)--(C);
   \node[site] at (A) {\(A\)}; \node[site] at (B) {\(B\)};
   \node[site] at (C) {\(C\)};
   \node[fill=white,text=rust,inner sep=2pt] at (O) {\(\ell\)};
   \node at (0,-.85) {One label for a rank term};
  \end{scope}
\end{tikzpicture}
\caption{The three tensor factors correspond to the two inputs and the
output paired against a third matrix. On the left, \(A,B,C\) see
\((i,j),(j,k),(k,i)\). On the right, all three see the same label
\(\ell\in[r]\) and apply their own linear maps.}\label{fig:three-wires}
\end{figure}

The critical ratio occurs at three tips of a tripod with unit arms.
The distances between tips are all two. If the shared-label rate is
\(8/3\), the ratio is \(6/(8/3)=9/4\). Under the strict-saving
hypothesis such a tripod cannot occur. Part~\ref{part:spheres} exploits
this exclusion. If varying the parameters of a tensor were to move
its preferred site in a certain family of prices, the motion would
produce a homogeneous tripod. A rectangular multiplication decoder
first determines the limiting rate when one arm is long. The rigidity
of polynomial multiplication then carries this equality to the tripod
with three equal arms (Proposition~\ref{prop:tripod}). Ruling out that
motion lets us price an entire family of tensors in a single extension
of the sites. This synchronization is the geometric input to compression. Once
the lower bound is proved, we can also identify the three wire lengths
at every price maximizing the value of \(T_n\)
(Corollary~\ref{cor:extremal-price}).

\subsection{From prices to a smaller network}

The family we use consists of sphere tensors. Their entries are
determined by marked points. When these points collide, a tensor
expands into contractions through legs of controlled dimension. Under
the strict-saving hypothesis we can also determine the largest prices
these tensors attain. For prescribed leg widths (logarithms of
dimensions), they attain asymptotically the cost of preparing a
tensor at the cheapest common site and transporting each leg to its
destination (Section~\ref{sec:sphere}). The marked points that specify
the entries and the sites that determine prices play different roles.
Figure~\ref{fig:org} shows how this calculation enters the two parts of
the network construction.

\begin{figure}[!htbp]
\centering
\begin{tikzpicture}[x=1cm,y=1cm,font=\small,
 box/.style={draw=ink,align=center,rounded corners=1.5pt,inner sep=5pt},
 arr/.style={draw=sea,thick,-{Latex[length=4pt]}}]
 \node[box] (price) at (0,0) {A strict square saving forces sphere prices\\
                                  {\footnotesize Part~\ref{part:spheres}}};
 \node[box,text width=4.0cm] (sort) at (-2.65,-1.65)
     {Stable sorting and\\batched table lookup\\{\footnotesize Part~\ref{part:sorting}}};
 \node[box,text width=4.0cm] (conc) at (2.65,-1.65)
     {Concentration of\\the incident bits\\{\footnotesize Part~\ref{part:local}}};
 \node[box] (vol) at (0,-3.5) {Compression of tensor powers\\
                                  {\footnotesize Part~\ref{part:volume}}};
 \draw[arr] (price.south west)--(sort.north);
 \draw[arr] (price.south east)--(conc.north);
 \draw[arr] (sort.south)--(vol.north west);
 \draw[arr] (conc.south)--(vol.north east);
\end{tikzpicture}
\caption{The implication needed for the lower bound. Sorting and
concentration give complementary estimates on the same network.
Repeated compression would violate the independent-list lower bound
of Theorem~\ref{thm:many}.}\label{fig:org}
\end{figure}

The compression must pay for each table of a network and for the
records that say how tables are connected. Direct lookup in a table
of \(u\) incident bits can cost \(2^u\). In a large batch, sorting
queries together with table entries amortizes that cost. Part~\ref{part:sorting}
uses incidence inequalities from sphere tensors to route records.
For \(2^u\) records of width \(O(u)\) the resulting stable sort costs
\(O(u^p)\) per record for some \(p<2\). Wider records admit a separate
estimate that keeps the cost close to linear in their width.

The strict inequality \(p<2\) matters because the quantity controlled
by collisions is proportional to the \emph{square} of a leg width.
Part~\ref{part:local} realizes the wiring with sphere tensors, moves
insertions across their seams and merges them. The use of fermionic contractions to reduce wiring is
related to the crossing construction of~\cite{Khoruzhii,oaG} and to
compression in subexponential satisfiability algorithms~\cite{falseeth,Anthropic}.
Merging concentrates the
incident bits at fewer stations while controlling the sum of squared
widths. Part~\ref{part:volume} applies these identities in rounds.
The expansion has many summands with different wirings. A single
network must describe the possible wirings, route the records, look
up the entries and sum the contributions. Counting all four operations
gives the logarithmic saving in Theorem~\ref{thm:shaving} and completes
the contradiction.

\subsection{Earlier bounds for matrix multiplication}\label{sec:history}

Strassen applied his seven-product identity recursively to blocks, giving
exponent $\log_2 7$~\cite{strassen}. Subsequent algorithms found savings
by working with several matrix products at once. Pan aggregated trilinear
forms~\cite{pan}, Bini showed how interpolation converts approximate
bilinear identities into exact algorithms~\cite{bini}, and Sch\"onhage's
asymptotic sum inequality turned a direct sum of matrix products into an
upper bound on the exponent of a single product~\cite{schoenhage}.
Coppersmith and Winograd studied recursive extraction of disjoint
products~\cite{cw82}. Strassen's laser method extracts them from tensor
powers~\cite{Strassen87}, and Coppersmith and Winograd designed an auxiliary
tensor for which progression-free sets accomplish the extraction~\cite{cw}.

The successive improvements depended on how many products could be kept
independent. Analyses of higher tensor powers~\cite{sto,ds,vw,legall},
refined extraction and asymmetric hashing~\cite{avw,duan,alphaomega,alman},
and numerical optimization~\cite{dupont} sharpened the bounds obtained
from the Coppersmith--Winograd tensor. Staggered extraction gave
$\omega_F<2.371055$ over every field~\cite{stagger}. A separate line of
work embedded matrix products in group algebras~\cite{cohn2003,cohn2005}
and in the algebras of coherent configurations~\cite{CohnUmansSupport}.

The upper bound $9/4$ comes from polynomial multiplication. In Strassen's
asymptotic spectrum~\cite{spectrum}, asymptotic rank is the maximum of
normalized functions that are additive on direct sums, multiplicative on
tensor products and monotone under restriction. Over $\mathbb C$, the
upper-bound proof tracks how these functions value polynomial multiplication
as the two input degrees grow. Degenerations give a lower bound on the
growth of this profile. Interpolation gives an upper bound that depends on
the value at matrix multiplication. Comparing the two bounds yields
$9/4$~\cite[Theorem~1.1]{nine}. The central degeneration separates blocks
that have independent variables on two legs and share the third. Partial
labels and semidirect sums give related formulations of this step
\cite{two258,li}. L\"u proved the separation step in arbitrary
characteristic~\cite{lu}. Amanov, van den Berg and Zuiddam obtained the
bound over every field~\cite[Theorem~1.1]{abvz}. See also the tensor-value
proof~\cite{tensorvalue} and Shpilka's explicit algorithms~\cite{shpilka}.
The rectangular version of the argument~\cite{uznanski,abvz} connects
these estimates to the dual exponent studied in~\cite{coppersmith,alphaomega}.

On the lower-bound side, Bl\"aser proved that the ordinary rank of $T_n$
is at least $\tfrac52 n^2-3n$ over every field~\cite{Blaeser}.
Over $\mathbb C$, Landsberg gave an ordinary-rank bound of
$3n^2-o(n^2)$~\cite{LandsbergRank}, and Landsberg and Micha{\l}ek
gave a border-rank bound of $2n^2-\lceil\log_2 n\rceil-1$
\cite{LandsbergMichalek}. A constant factor above $n^2$ disappears on
taking roots along the sequence $T_n^{\otimes j}\simeq T_{n^j}$.
Barriers for the laser method and bounds from tensor irreversibility
address algorithms that pass through specified intermediate
tensors~\cite{limits,irreversible}. For polynomial multiplication,
Amanov, van den Berg and Zuiddam similarly proved that $9/4$ is optimal
for a specified set of degenerations and combination
rules~\cite[Theorem~5.4]{abvz}. Theorem~\ref{thm:square} gives the
matching lower bound for unrestricted arithmetic circuits. Related
work on the spectrum constructs quantum functionals~\cite{CVZ} and
explicit points over arbitrary fields~\cite{almanlipratt}, and describes
the sublevel sets of asymptotic rank by polynomial
equations~\cite{SpectralSublevels}. Wigderson and Zuiddam survey this
framework~\cite{WZ}.

